% Лекц 4: семинар ба гэрийн даалгавар. Бодлогууд: Г.Н. Берман, "Сборник задач по курсу
% математического анализа" (22-р хэвлэл). [Б.N] нь номын бодлогын дугаар.

\ProblemSet{Семинар 4}
\begin{frame}[allowframebreaks]{Семинар 4}
\pgroup{Функцийн тодорхойлогдох мужийг ол.}
\begin{plist}
\prob{47.3}{$y=\sqrt{5-2x}$}
\prob{47.8}{$y=\dfrac{2x}{x^2-3x+2}$}
\prob{47.11}{$y=\sqrt{x^2-4x+3}$}
\prob{47.12}{$y=\dfrac{x}{\sqrt{x^2-3x+2}}$}
\prob{47.18}{$y=\sqrt{1-|x|}$}
\prob{48.12}{$y=\sqrt{x^2-3x+2}+\dfrac{1}{\sqrt{3+2x-x^2}}$}
\prob{48.13}{$y=(x^2+x+1)^{-3/2}$}
\end{plist}
\pgroup{Функцийн утга, тэгш ба сондгой функц.}
\begin{plist}
\prob{107}{$f(x-1)=2x^2-3x+1$ бол $f(x+1)$-ийг ол.}
\prob{54.1}{$y=x^4-2x^2$ тэгш, сондгой, аль нь ч биш үү?}
\prob{54.2}{$y=x-x^2$ тэгш, сондгой, аль нь ч биш үү?}
\prob{54.5}{$y=x-\dfrac{x^3}{6}+\dfrac{x^5}{120}$ тэгш, сондгой, аль нь ч биш үү?}
\prob{55.1}{$y=x^2+3x+2$ функцийг тэгш ба сондгой функцийн нийлбэрээр илэрхийл.}
\end{plist}
\pgroup{Олон гишүүнт функцийн шинж ба график.}
\begin{plist}
\prob{53.4}{$y=x^3-3x^2+2x$ функцийн тэг ба тэмдэг тогтмол байх завсруудыг ол.}
\prob{77.1}{$y=|x+1|+|x-1|$ функцийн графикийг байгуул.}
\prob{81.9}{$y=2x^2-6x+4$ функцийн графикийг байгуулж, өсөх ба буурах завсрыг заа.}
\prob{83.2}{$y=-x^2-3x+2$ функцийн хамгийн их утгыг ол.}
\end{plist}
\pgroup{Урвуу функц ба бутархай-шугаман функц.}
\begin{plist}
\prob{117.3}{$y=1-3x$ функцийн урвуу функцийг ол.}
\prob{117.6}{$y=\dfrac{1}{1-x}$ функцийн урвуу функцийг ол.}
\prob{113.1}{$y=\dfrac{x-1}{x-2}$ функцийн графикийг байгуул.}
\prob{114.2}{$y=\dfrac{x}{2x-5}$ функцийн $[-1,2]$ хэрчим дээрх хамгийн их, бага утгыг графикаас ол.}
\end{plist}
\end{frame}

\ProblemSet{Гэрийн даалгавар 4}
\begin{frame}[allowframebreaks]{Гэрийн даалгавар 4}
\pgroup{Функцийн тодорхойлогдох мужийг ол.}
\begin{plist}
\prob{47.5}{$y=\dfrac{1}{x^2-1}$}
\prob{47.7}{$y=\dfrac{1}{x^3-x}$}
\prob{47.9}{$y=1-\sqrt{1-x^2}$}
\prob{47.10}{$y=\dfrac{1}{\sqrt{x^2-4x}}$}
\prob{47.19}{$y=\dfrac{1}{\sqrt{|x|-x}}$}
\prob{47.20}{$y=\dfrac{1}{\sqrt{x-|x|}}$}
\prob{48.5}{$y=\sqrt{x-1}+2\sqrt{1-x}+\sqrt{x^2+1}$}
\prob{48.11}{$y=\sqrt{\dfrac{x-2}{x+2}}+\sqrt{\dfrac{1-x}{\sqrt{1+x}}}$}
\end{plist}
\pgroup{Функцүүд адилхан уу?}
\begin{plist}
\prob{49.1}{$f(x)=\dfrac{x}{x^2}$ ба $\varphi(x)=\dfrac1x$}
\prob{49.2}{$f(x)=\dfrac{x^2}{x}$ ба $\varphi(x)=x$}
\prob{49.3}{$f(x)=x$ ба $\varphi(x)=\sqrt{x^2}$}
\end{plist}
\pgroup{Тэгш ба сондгой функц, тэмдэг.}
\begin{plist}
\prob{53.2}{$y=x^2-5x+6$ функцийн тэг ба тэмдэг тогтмол байх завсруудыг ол.}
\prob{54.8}{$y=1-x^2$ тэгш, сондгой, аль нь ч биш үү?}
\prob{55.2}{$y=1-x^3-x^4-2x^5$ функцийг тэгш ба сондгой функцийн нийлбэрээр илэрхийл.}
\prob{56}{$f(x)+f(-x)$ тэгш, $f(x)-f(-x)$ сондгой функц болохыг батал.}
\prob{58}{Хоёр тэгш функцийн үржвэр тэгш, хоёр сондгой функцийн үржвэр тэгш, тэгш ба сондгой функцийн үржвэр сондгой функц болохыг батал.}
\end{plist}
\pgroup{График, хамгийн их ба бага утга.}
\begin{plist}
\prob{76}{$\varphi(x)=\frac x2+2$ ($x\le2$), $\varphi(x)=5-x$ ($x\ge2$) бол $\varphi(x)=2x-4$ тэгшитгэлийг аналитик ба графикийн аргаар бод.}
\prob{77.2}{$y=|x+1|-|x-1|$ функцийн графикийг байгуул.}
\prob{77.3}{$y=|x-3|-2|x+1|+2|x|-x+1$ функцийн графикийг байгуул.}
\prob{81.3}{$y=|x^2-1|$ функцийн графикийг байгуулж, өсөх ба буурах завсрыг заа.}
\prob{81.12}{$y=-x|x|$ функцийн графикийг байгуулж, өсөх ба буурах завсрыг заа.}
\prob{84.3}{$y=1-3x+6x^2$ функцийн хамгийн бага утгыг ол.}
\prob{85}{$a$ тоог үржвэр нь хамгийн их байх хоёр нэмэгдэхүүний нийлбэрт задал.}
\prob{87}{Чулуун хананы дэргэд нийт 8 м урт модон хашаагаар тэгш өнцөгт талбай хашина. Талбай хамгийн их байхын тулд хашааны хананд параллель хэсэг ямар урт байх вэ?}
\end{plist}
\pgroup{Урвуу ба бутархай-шугаман функц.}
\begin{plist}
\prob{117.4}{$y=x^2+1$ функцийн урвуу функцийг ол.}
\prob{117.7}{$y=x^2-2x$ функцийн урвуу функцийг ол.}
\prob{118}{$y=\dfrac{ax+b}{cx+d}$ ($ad-bc\ne0$) бутархай-шугаман функцийн урвуу функц мөн бутархай-шугаман болохыг батал.}
\prob{119}{Ямар нөхцөлд бутархай-шугаман функц өөрийнхөө урвуутай давхцах вэ?}
\prob{113.2}{$y=\dfrac{2x}{3-x}$ функцийн графикийг байгуул.}
\prob{116}{Графикуудыг нэмэх аргаар $y=\dfrac{x^2+1}{x}$ функцийн графикийг байгуул.}
\end{plist}
\end{frame}
